\documentclass[math, info]{mpb-cours}

\title{Modèles Déformables et Chemins Minimaux en Analyse d'Images\texorpdfstring{\\ Rapport sur \textit{Jet Marching Methods for Solving the Eikonal Equation}}{}}
\author{Matthieu Boyer}


\begin{document}
\bettertitle
\newpage
\section{Étude Synthétique}

\subsection{Problème Traité}
Dans \cite{potter2021jetmarchingmethodssolving}, les auteurs s'intéressent à la résolution de l'équation
eikonale.
Leur objectif est notamment d'améliorer la précision de la méthode de \textit{Fast Marching} afin de
résoudre des problèmes en acoustique haute-fréquence.

\smallskip

Les méthodes usuelles de résolution de ce genre de problèmes sont basées sont sur des méthodes de type
ray-tracing qui demandent de résoudre avec une bonne précision l'équation de transport de l'amplitude.
Pour ce faire, il est nécessaire de calculer les dérivées partielles de deuxième ordre de l'eikonale.
S'il existe des solveurs\footnote{Faute de meilleur mot, on utilisera ici l'anglicisme "solveur" pour
	\textit{numerical solver}} pour l'équation eikonale, leur qualité est réduite lorsqu'on considère un
maillage approximatif de l'espace de propagation ou même simplement un espace de propagation d'une forme
complexe.

Les méthodes basées sur le raytracing sont en effet vulnérables aux effets au bord des caustiques, les
enveloppes des faisceaux réfléchis et réfractés, et les fortes variations de la vitesse de l'onde
engendrent de nombreuses caustiques distinctes dans le cas par exemple des ondes sismiques.
Ensuite, bien que la vitesse varie très lentement dans une pièce, les caustiques au bord des objets dans une
pièce complexifient fortement le calcul.

\subsection{Équations et Méthodes Numériques}
On s'intéresse à la résolution en $\tau$ de l'équation Eikonale sur un domaine $\Omega \subseteq \R^{n}$ avec
un bord $\partial \Omega$ et une condition au bord $\Gamma \subseteq \Omega$
\begin{equation*}
	\begin{cases}
		\norm{\nabla \tau(x)} = s(x) & x \in \Omega \\
		\tau(x) = g(x)               & x\in \Gamma
	\end{cases}.
\end{equation*}
L'équation eikonale apparaît notamment dans le cadre de la recherche des modes propres de l'équation de
D'Alembert avec l'ansatz BKW\footnote{Forme générale d'une solution qui apparaît notamment dans l'étude de
	l'équation de Schrödinger} qui est de la forme $\alpha(x)e^{i\omega \tau(x)}$.

\smallskip

En se plaçant de plus dans le cadre de l'optique géométrique (à haute fréquence), $\tau$ sera une solution
de l'équation eikonale tandis que $\alpha$, l'amplitude, vérifiera l'équation de transport
\begin{equation*}
	\alpha\Delta\tau + 2\nabla\tau^{T}\nabla\alpha = 0.
\end{equation*}
$\tau$ est alors donné par le principe de Fermat\footnote{"\textit{Le chemin suivi par un rayon entre deux
		points est celui minimisant le temps.}", il découle du principe de Maupertuis "\textit{Le chemin suivi
		par une particule en mouvement minimise son action.}"}
\begin{equation*}
	\tau(x) = \min_{y\in \Gamma, \psi: x \rightsquigarrow y} \tau(y) + \int_{0}^{1} s(\psi(\sigma))\norm{\psi'(\sigma)}\d\sigma.
\end{equation*}

Par la récursion présente dans la formulation du principe de Fermat, on va se baser une méthode de
propagation, à la manière du Fast Marching de Sethian.

À la place de considérer $\Gamma$ et $\Omega$ comme des ensembles dans $\R^{n}$, on demandera qu'ils soient
des maillages\footnote{Dans l'implémentation que nous étudierons plus bas, il s'agira même de parties de
	triangulations de surfaces plongées dans $\R^{3}$.} de taille caractéristique $h$, et on va propager
$T : \Omega_{h} \to \R$ une approximation numérique de l'eikonale ainsi que son gradient.
On appliquera ensuite une version modifiée de l'algorithme de Dijkstra permettant de calculer sur chaque
face du maillage une approximation cubique du rayon permettant de propager le résultat, grace à une approche
semi lagrangienne.

\subsection{Rapport au Cours}
La méthode proposée par les auteurs est une extension directe de l'algorithme de Fast Marching,
utilisant le \textit{jet} (l'ensemble des dérivées d'ordre un) de l'eikonale.
Cet article propose donc une méthode efficace (dépendamment de la précision souhaitée~!) pour résoudre
précisément l'équation eikonale, bien que tournée à la base vers l'étude de problèmes acoustiques.

\subsection{Originalité du Travail}
Le travail en lui-même ne développe pas un "nouvel algorithme", mais améliore fortement l'algorithme de
Fast Marching en précision, et dépasse le schéma eulérien parmi les solveurs de l'équation Eikonale.
L'originalité de cette méthode se trouve dans l'intégration d'un schéma différentiels au sein d'un algorithme
propagatif, tout en restant optimaleement local et flexible.

\subsection{Résultats Principaux}
L'algorithme développé dans l'article permet une précision $\O(h^{3})$ en la taille de l'approximation
du maillage $h$ sur la précision de $T$ et en un ordre de moins pour la précision du gradient, tout en
restant en temps $\O(N\log(N)\times P)$, où $P$ est le temps de résolution d'un problème de minimisation
sur une face, qui est borné par $N$ mais est réellement de l'ordre du degré moyen d'un sommet dans le
maillage à la puissance de la dimension (dans notre cas, rarement plus que $6^{2}$).

\subsection{Faiblesses}
Le principal défaut de la méthode proposé dans \cite{potter2021jetmarchingmethodssolving} est la complexité
de celle-ci en haute dimension.
En effet, le nombre de problèmes de minimisation à résoudre est exponentiel en la dimension, et leur
complexité est polynomiale en celle-ci.
Toutefois, ces problèmes sont aisément contournables lorsqu'on a un peu plus d'informations sur la structure
de l'espace de propagation.

Ensuite, le second problème vient avec la complexité liée au nombre de possibilités pour l'algorithme.
Celui-ci permet de nombreuses variations différentes et est donc difficilement adaptable pour un
problème spécifique, demandant un grand temps d'adaptation.
Ceci est cependant un faux problème qui ne vient pas vraiment de l'algorithme en lui-même.

Enfin, l'algorithme est numériquement assez dépendant de la qualité du maillage: pour un obtenir une
précision numérique $\epsilon = h^{3}$, il faut de l'ordre $\lambda_{d}(\Omega) d/ h^{d}$ points pour
discrétiser $\Omega$, où $\lambda_{d}$ est la mesure de Lebesgue en $d$ dimensions.
Ceci n'est pas forcément un problème en pratique, lorsqu'on se limite à des surfaces.

\newpage

\section{Description et Résumé de l'Article}

Dans ce rapport sur \textit{Jet Marching Methods for Solving the Eikonal Equation} par Potter et Cameron,
nous étudierons leur extension de l'algorithme de \textit{Fast Marching} utilisant les dérivées au
premier ordre pour améliorer la précision de solveurs de l'équation eikonale.

\subsection{Formulation du problème}
On s'intéresse à la recherche de solutions de l'équation d'onde de Helmholtz, sous la forme de l'ansatz
BKW de la mécanique quantique $u(x) \sim \alpha(x)e^{i\omega\tau(x)}$.
Ici, on connaîtra la vitesse $s$ de l'onde sur $\Omega \subseteq \R^{n}$, et son temps de départ $g$ sur
$\Gamma$.
De l'équation de Helmholtz, on en tire deux équations principales:
\begin{enumerate}
	\item L'équation \define{eikonale} vérifiée par $\tau$.
	      \begin{equation}
		      \begin{cases}
			      \norm{\nabla \tau(x)} = s(x) & x \in \Omega \\
			      \tau(x) = g(x)               & x\in \Gamma
		      \end{cases}. \label{eq:eik}
	      \end{equation}
	      En particulier, $\tau$ est définie par le principe de Fermat
	      \begin{equation}
		      \tau(x) = \min_{y\in \Gamma, \psi: x \rightsquigarrow y} \tau(y) + \int_{0}^{1} s(\psi(\sigma))\norm{\psi'(\sigma)}\d\sigma.\label{eq:fermat}
	      \end{equation}
	\item L'équation de \define{transport} de l'amplitude
	      \begin{equation}
		      \alpha\Delta\tau + 2\nabla\tau^{T}\nabla\alpha = 0.\label{eq:amp}
	      \end{equation}
\end{enumerate}

\smallskip

On s'intéressera en pratique à la résolution du problème ci-dessus sur un maillage $\Omega_{h}$ de taille
caractéristique $h$ d'une variété de dimension $2$ plongée dans $\R^{3}$.
Dans ce cadre, on cherchera à calculer $T$ et $A$ des approximations de $\tau$ et $\alpha$ sur le maillage
avec une bonne précision, de manière localement optimale et dérivable.

\subsection{L'algorithme de Jet Marching}
L'équation du principe de Fermat \eqref{eq:fermat} permet de définir la solution $\tau$ de l'équation
Eikonale \eqref{eq:eik} à partir d'un ensemble fini de points.
En particulier, on peut utiliser une approche de propagation de solutions basée sur le principe d'optimalité
de Bellman pour calculer une solution générale.

\smallskip

À la manière du \textit{Fast Marching} et de l'algorithme de Dijkstra, sur le maillage $\Omega_{h}$ vu comme
un graphe, à partir d'un ensemble de sources $\Gamma_{h}$.
La principale différence dans l'algorithme de \define{jet marching} découle de ce qui va être propagé
comme approximation numérique: fast marching ne propage que le temps d'arrêt local, tandis que le jet
marching va aussi propager le gradient de l'approximation numérique: le 1-jet de l'approximation.
Ensuite, une fois qu'on sait que notre approximation ne sera pas améliorée en un point (les sources étant
fixées car des minima), on peut résoudre localement l'équation du principe de
Fermat \eqref{eq:fermat} pour obtenir l'approximation sur de nouveaux points et répéter l'opération jusqu'à
épuiser les points sans valeur.

\smallskip

Donnons une description haut niveau un peu plus formelle de l'algorithme:
On va attribuer à chaque sommet une étiquette parmi \texttt{valide}, \texttt{jugé}, et
\texttt{loin}.
On initialise d'abord l'approximation sur les sources dans $\Gamma_{h}$, qu'on étiquette \texttt{valide}.
On initialise aussi une file de priorité avec les voisins des sources, qu'on étiquette \texttt{jugé}.
Tant que la pile n'est pas vide, on dépile le minimum (par distance de l'approximation en cours), puis,
si le sommet n'est pas \texttt{valide}, on évalue une nouvelle approximation du $1$-jet de $T$,
on ajoute ses voisins à la pile en les marquant \texttt{jugé}, et on le marque \texttt{valide}.
On résout alors le calcul de son amplitude à partir de la trajectoire approximée du rayon.

\smallskip

Cet algorithme est donc très similaire à la méthode de fast marching, en remplaçant simplement une
approximation numérique par une approximation du $1$-jet.
Il y a toutefois quelques points de détail à préciser:
\begin{enumerate}
	\item Comment représente-t-on le rayon portant l'eikonale localement~?
	\item Comment propage-t-on l'approximation numérique du $1$-jet~?
	\item Comment approxime-t-on l'amplitude~?
\end{enumerate}
On va présenter les solutions possibles à ces questions ci-dessous.

\subsubsection{Propagation du \texorpdfstring{$1$-Jet}{1-Jet}}
Afin de propager le $1$-jet de $T$ en un point objectif $\hat{x}$, on va considérer un ensemble de
voisins $(\hat{x}, x_{1}, \ldots, x_{d})$ de $\hat{x}$ étiquetés \texttt{valide}.
On appellera un tel ensemble un \define{pochoir de mise à jour} et on appellera $d$ sa \define{dimension}.
L'idée va être de résoudre l'équation Eikonale localement via la minimisation d'une fonctionnelle basée
sur le principe de Fermat \eqref{eq:fermat}.

\smallskip

Ici, il est important de noter que $d$ ne peut pas être plus grand que la dimension de $\Omega$.
On va adopter une méthode de bas en haut pour calculer le pochoir donnant la mise à jour optimale:
on regarde le pochoir de dimension $1$ (donc le voisin $x_{1}$) donnant la mise à jour optimale, puis on
fixe $x_{1}$ et on regarde les pochoirs de dimension $2$ contenant $x_{1}$ pour trouver $x_{2}$
donnant la mise à jour optimale, et ainsi de suite.

\smallskip

Pour chaque pochoir, on va chercher le point $x_{\lambda}$ dans le $(d - 1)$-simplexe
$\mathrm{conv} \{x_{1}, \ldots, x_{d}\}$, de sorte que la courbe $\psi$ joignant $x_{\lambda}$ à $\hat{x}$
selon \eqref{eq:fermat} donne une valeur de $\tau$ en $\hat{x}$ qui soit minimale.
Il faut donc résoudre un problème de minimisation par $x_{\lambda}$ pour calculer la courbe $\psi$, et un
problème de minimisation pour trouver le $x_{\lambda}$ adapté.
Heureusement pour nous, le problème de minimisation pour $x_{\lambda}$ est simple: comme nous le verrons
plus bas, les fonctionnelles de Fermat donnent des problèmes strictement convexes en $x_{\lambda}$ sur
un ensemble de définition convexe, donc en particulier, ils sont résolubles avec une bonne précision en
temps constant par la méthode de Newton par exemple.
Il nous suffit donc simplement de savoir calculer la courbe $\psi$ associée à $x_{\lambda}$ pour pouvoir
terminer l'algorithme.


\subsubsection{Principe de Fermat}
Résoudre la mise à jour associée à un pochoir et un point $x_{\lambda}$ donnés demande d'abord de savoir
représenter la courbe $\psi$ solution de l'équation du principe de Fermat.
Dans le cas où la vitesse de l'onde dans le milieu est constante, il s'agit d'une droite, mais
dans le cas où elle ne l'est pas, la fonction peut-être arbitrairement compliquée.

\smallskip

Pour simplifier, on va utiliser une approximation cubique $\phi$ de la trajectoire du rayon.
Notons $L = \norm{\hat{x} - x_{\lambda}}$, de sorte que $\psi: [0, L] \to \Omega$.
On se donne deux paramètres $t_{\lambda}$ et $\hat{t}$ de sorte que $\phi$ soit définie par
\begin{equation}
	\phi(0) = x_{\lambda} \quad \phi'(0) \sim t_{\lambda} \quad \phi(L) = \hat{x} \quad \phi'(L) \sim \hat{t}.
\end{equation}
La fonctionnelle de coût de Fermat (l'intégrale de \eqref{eq:fermat}) s'écrit alors
\begin{equation}
	F(\phi) = T(x_{\lambda}) + \frac{L}{6}\left(s(x_{\lambda}\norm{\phi'(0)}) + 4s(\phi_{1/2})\norm{\phi_{1/2}'} + s(\hat{x})\norm{\phi'(L)}\right),
\end{equation}
où $\phi_{1/2} = \phi(L/2)$ et $\phi_{1/2}' = \phi'(L/2)$.

\smallskip

On a alors plusieurs possibilités pour la représentation de $\phi$ et des quatre paramètres nécessaires:
\begin{enumerate}
	\item Via une perturbation $\delta\phi$ à partir de la ligne droite $l$ entre $x_{\lambda}$ et
	      $\hat{x}$: ici, $\delta\phi$ est définie par deux vecteurs unitaires $t_{\lambda}$ et $\hat{t}$
	      donnant la vitesse en $0$ et en $L$;
	\item Via le graphe d'une fonction $\zeta$ orthogonale à $l'$: si on considère $Q$ orthogonale telle que
	      $Q^{T}l' = 0$, on choisit $\delta\phi = Q\zeta$ et $\zeta$ est paramétrisée par les dérivées de la
	      perturbation.
\end{enumerate}
Les deux manières sont équivalentes mais donnent des problèmes de minimisation un peu différents.
On ne rentrera pas ici dans le détail des différences, mais il est impartant de noter que la précision
de l'interpolante dépend assez fortement de la qualité de notre approximation.

\smallskip

Il faut ensuite noter que connaître $\hat{t}$ va nous permettre de propager le gradient de $T$ en plus de
sa valeur.
En effet, $\psi$ étant une solution de l'intégrale de Fermat, c'est une caractéristique de l'eikonale et
son vecteur tangent doit être localement parallèle à $\nabla\tau$ et donc, après minimisation,
on peut simplement assigner $s(\hat{x})\hat{t}$ à $\nabla T(\hat{x})$ ce qui propage effectivement
le $1$-jet.

\smallskip

Finalement, il reste simplement à savoir comment calculer la caractéristique de l'eikonale.
\cite{potter2021jetmarchingmethodssolving} proposent quatre méthodes différentes, chacune adaptable à
la représentation choisie pour l'approximante $\phi$:
\begin{enumerate}
	\item Par minimisation directe: cette méthode est la plus simple, on optimise directement le problème en
	      $x_{\lambda}$, et $t_{\lambda}, \hat{t}$ unitaires. C'est un problème de dimension
	      $(d - 1)(n - 1)^{2}$.
	\item Par l'eikonale: on utilise l'équation Eikonale sur le simplexe pour retrouver une relation
	      permettant de supprimer le paramètre $t_{\lambda}$ via les dérivées directionnelles sur le simplexe.
	      C'est un problème de dimension $(d - 1)(n - 1)$.
	\item Par une marche d'interpolants sur les cellules: on simplifie la méthode précédente en retirant
	      l'utilisation de dérivées directionnelles et en utilisant la connaissance d'une cellule précédente.
	      C'est un problème de dimension $(d - 1)(n - 1)$, mais qui permet de marcher les dérivées
	      partielles de second ordre de $T$.
	\item Dans des cas particuliers, par exemple quand la vitesse du son est linéaire, on peut se
	      restreindre à une interpolante quadratique, ce qui retire un paramètre.
	      C'est un problème de dimension $(d - 1)(n - 1)$.
\end{enumerate}

\subsubsection{Évaluation de l'Amplitude}
Pour évaluer l'amplitude numérique $A$ en $\hat{x}$ à partir de $T$ et de $\nabla T$, une idée pourrait
être de discrétiser l'équation de transport \eqref{eq:amp}.
Cependant, ceci causerait l'approximation d'être singulière aux caustiques.
Pour éviter ce problème, on utilise une méthode tirée du raytracing paraxial: à partir d'un rayon
$\phi_{0}$ (celui généré par la mise à jour) on considère des rayons dans un tube autour.

On peut alors mesurer la divergence géométrique autour du rayon central d'une solution, et en déduire une
valeur numérique pour l'approximation $A(\hat{x})$.
Nous ne rentrerons pas dans les détails de cette partie, mais \cite{potter2021jetmarchingmethodssolving}
propose de construire une solution des équations d'Euler-Lagrange pour l'équation Eikonale le long de
$\phi$ avec une position $Q$ et un moment $P$ et de calculer $\abs{\det Q(L)}$ pour estimer la
divergence du faisceau de rayons.

Alors, on pourra écrire
\begin{equation}
	A(\hat{x}) = \sqrt{\frac{\frac{1}{s(\phi_{0}(L))}}{\abs{\det Q(L)}}}\mathit{A}(x_{\lambda}),
\end{equation}
ce qui nous donne une approximation de $A(\hat{x})$ à partir d'une approximation polynomiale $\mathit{A}$
de $A$ basée sur les valeurs de $A$ sur $\Omega_{h}$.

\newpage

\section{Implémentation}
Dans cette section, je présente mon implémentation du jet marching et ses résultats.

\subsection{Déclaration d'honnêteté}
Pour le cours de \textit{Geometric Data Analysis}, j'avais avec Antoine Groudiev écrit une application
en Rust sur le calcul par diverses méthodes des géodésiques.
Nous nous étions concentrés sur des $2$-variétés plongées dans $\R^{3}$ afin de pouvoir visualiser.
Dans le cadre de ce projet précédent, j'avais personnellement implémenté des méthodes basées sur des
équations différentielles tirées de la physique et notamment l'algorithme de Fast Marching (basé, donc,
sur l'équation Eikonale).

\smallskip

Dans le cadre de ce projet, j'ai réutilisé la structure de données représentant en mémoire les variétés,
celle utilisée pour les problèmes géodésiques (qui représente les conditions initiales),
ainsi que le visualisateur 3D écrit par Antoine.
J'ai partiellement repris la structure de l'algorithme de fast marching, mais en réimplémentant tout de
zéro.

Le code est disponible, avec l'historique du code du précédent projet sur ce dépôt git: \url{pouet}

\subsection{Détails d'implémentation}
Dans cette implémentation, contrairement à la méthode générale décrite par Potter et al. dans
\cite{potter2021jetmarchingmethodssolving}, je me suis limité à des $2$-variétés connexes plongées
dans $\R^{3}$.
Les surfaces en questions sont représentées par un maillage triangulaire: les ensembles $\Omega_{h}$ et
$\Gamma_{h}$ sont définis à partir d'un ensemble $\mathcal{V}$ de sommets et un ensemble
$\mathcal{F}\subseteq \mathcal{V}^{3}$ de faces triangulaires.

Les pochoirs de mise à jour au c\oe ur de l'algorithme sont donc limités à la dimension au plus $2$, et les
sommets à choisir dans l'ensemble de voisinage sont choisis d'une manière parmi deux une fois que le premier
sommet $x_{1}$ est fixé par les mises à jour d'arête vers $\hat{x}$:
\begin{enumerate}
	\item où bien on considère $x_{2}$ tel que $(\hat{x}, x_{1}, x_{2}) \in \mathcal{F}$;
	\item où bien on considère $x_{2}$ tel que $\max(\norm{\hat{x} - x_{2}}_{1}, \norm{x_{1}, x_{2}}_{1})$
	      soit suffisamment petit.
\end{enumerate}

\subsection{Résultats}

\bibliographystyle{alpha}
\bibliography{report}
\end{document}
